\documentclass[11pt,a4paper]{article}
\usepackage[a4paper,margin=2.4cm]{geometry}
\usepackage{graphicx}
\usepackage{float}

\title{Labwork 3}
\author{Pham Tien Nhat \\ Student ID: 2540090 \\ Advanced HPC -- USTH}
\date{\today}

\begin{document}
\maketitle

\section{Loading the Image}
\texttt{pixelSum} $= h \times w = 306240$.

\begin{verbatim}
import matplotlib.pyplot as plt

image = plt.imread('../../image/images.jpeg')
plt.imshow(image)
h ,w ,d = image.shape

pixelSum = h * w
\end{verbatim}

\section{Flatten}
The image is flattened into a 1D array of RGB pixels.

\begin{verbatim}
flatten = image.reshape(pixelSum,3)
\end{verbatim}

\section{Grayscale on CPU}
A \texttt{for} loop computes $g = (R+G+B)/3$ per pixel.

\begin{verbatim}
import numpy as np
import time 
grayCPU = np.zeros(pixelSum, np.uint8)
start = time.time()
for i in range(pixelSum):
    grayCPU[i] = np.uint8((int(flatten[i, 0]) + int(flatten[i, 1])
                          + int(flatten[i, 2])) / 3)
end = time.time() - start
print(end)
plt.imshow(grayCPU.reshape(h, w), cmap='gray')
\end{verbatim}

\section{Grayscale on GPU}

\begin{verbatim}
@cuda.jit
def grayscale(src, dst):
    tidx = cuda.threadIdx.x + cuda.blockIdx.x * cuda.blockDim.x
    if tidx >= src.shape[0]:
        return
    g = np.uint8((src[tidx, 0] + src[tidx, 1] + src[tidx, 2]) / 3)
    dst[tidx, 0] = dst[tidx, 1] = dst[tidx, 2] = g
\end{verbatim}

The host copies data from GPU to the device, launches the kernel and copies the result
back.

\begin{verbatim}
blockSize = 64
gridSize = (pixelSum + blockSize - 1) // blockSize
start_gpu = time.time()
devSrc = cuda.to_device(flatten)
devDst = cuda.device_array((pixelSum, 3), np.uint8)
grayscale[gridSize, blockSize](devSrc, devDst)
hostDst = devDst.copy_to_host()
end_gpu = time.time() - start_gpu
print(end_gpu)
plt.imshow(hostDst.reshape(h, w, 3))
\end{verbatim}

\section{Block Size vs Time}
The kernel is timed for block sizes $2^5$ to $2^{10}$.

\begin{verbatim}
blockSizes = [32, 64, 128, 256, 512, 1024]
times = []
grayscale[(pixelSum + 63) // 64, 64](devSrc, devDst) 
for blockSize in blockSizes:
    gridSize = (pixelSum + blockSize - 1) // blockSize
    start = time.time()
    for i in range(100):
        grayscale[gridSize, blockSize](devSrc, devDst)
    cuda.synchronize()
    times.append((time.time() - start) / 100)
plt.plot(blockSizes, times, marker='o')
plt.xscale('log', base=2)
plt.xlabel('block size')
plt.ylabel('time')
plt.show()
\end{verbatim}

\begin{figure}[H]
\centering
\includegraphics[width=0.7\textwidth]{../../image/image.png}
\caption{Block size vs kernel time (seconds).}
\end{figure}

\end{document}
